\section{What Selection Carries Forward}

This chapter has made relativity a grammatical selection rather than a view
from nowhere. A frame is the permitted quotient by which observationally
equivalent presentations are identified. Relative velocity is the residue
left after that quotient has done its work. The grammar therefore forbids
absolute motion as an unearned possession and permits relative motion as a
carried discrepancy between reconciled records.

The same selection discipline read charge. An open path can carry local
transport, but it does not close an orientation and therefore carries no
charge. A closed loop can return with a holonomy sign. The flat native
baseline reads the negative sign. A tilted or oriented closure may read the
positive sign. The positron is that positive sign of closure, not an
imported second stock of being.

The native asymmetry followed as a count. Over the symmetric native
baseline, the grammar counts matter and no antimatter. The zero antimatter
count is not ignorance; it is the verdict of a specified selection that
forbids the positive sign there. Antimatter is admitted only when the
comparison tilts or otherwise selects the asymmetric closed sign.

The finite gauge then housed the count. Its thirty-six gates partition into
pooling and selecting moves. Pooling preserves the native matter reading.
Selecting turns the asymmetric positive sign. The all-pooling corner reads
native matter with no antimatter. The one-crank pair admits a single
positive turn while the remaining gates continue to pool. Matter and
antimatter together exhaust the gauge only when every gate has been assigned
to one side of the partition.

Finally, the predicates turned back on the reader. Did this distinguish? Did
this fall within the magnitude band? Did this rise? Did this gate pool or
select? Did the loop close? The reader is the final station because those
questions have no further interior hiding place. But the reader is
constrained by the grammar that formed the questions. Selection is lawful
only where equivalence, orientation, count, and boundary have been earned.

What passes forward is a closed obligation. If relative velocity is quotient
selection, if charge is loop sign, if native asymmetry is a baseline count,
if the gauge partitions every gate, and if the reader is the final station,
then the next movement cannot simply add another forward structure. It must
walk the tower back. Chapter 08 receives a reader-facing loop and asks how
the construction returns through its own steps to the needle that first made
identity admissible.
